%%%%%%%%%%%%%%%%%%%%%%%%%%%%%%%%%%%%%%%%%%%%%%%%%%%%%%%%%%%%%%%%%%%%%%%%%%%%%%%%
% TUM-Vorlage: Plakat A0 Querformat
%%%%%%%%%%%%%%%%%%%%%%%%%%%%%%%%%%%%%%%%%%%%%%%%%%%%%%%%%%%%%%%%%%%%%%%%%%%%%%%%
%
% Rechteinhaber:
%     Technische Universität München
%     https://www.tum.de
% 
% Gestaltung:
%     ediundsepp Gestaltungsgesellschaft, München
%     http://www.ediundsepp.de
% 
% Technische Umsetzung:
%     eWorks GmbH, Frankfurt am Main
%     http://www.eworks.de
%
%%%%%%%%%%%%%%%%%%%%%%%%%%%%%%%%%%%%%%%%%%%%%%%%%%%%%%%%%%%%%%%%%%%%%%%%%%%%%%%%

%%%%%%%%%%%%%%%%%%%%%%%%%%%%%%%%%%%%%%%%%%%%%%%%%%%%%%%%%%%%%%%%%%%%%%%%%%%%%%%%
%\documentclass[a0, plainsections, landscape]{sciposter}
\newcommand{\PlakatFormat}{A0 quer}

\documentclass{scrlttr2} % Dokumentenklasse: KOMA-Skript Letter

% Anpassbare Kopf- und Fußzeilen für KOMA-Skript:
%\usepackage{scrlayer-scrpage}

\usepackage[utf8]{inputenc} % Textkodierung: UTF-8
\usepackage[T1]{fontenc} % Zeichensatzkodierung

\usepackage{calc} % Berechnungen

\usepackage[ngerman]{babel} % Deutsche Lokalisierung
\usepackage{graphicx} % Grafiken
\usepackage[absolute]{textpos} % Positionierung

% Silbentrennung:
\usepackage{hyphenat}
%\tolerance 2414
%\hbadness 2414
%\emergencystretch 1.5em
%\hfuzz 0.3pt
%\widowpenalty=10000     % Hurenkinder
%\clubpenalty=10000      % Schusterjungen
%\vfuzz \hfuzz

% Euro-Symbol:
\usepackage[gen]{eurosym}
\DeclareUnicodeCharacter{20AC}{\euro{}}

% Schriftart Helvetica:
\usepackage[scaled]{helvet}
\renewcommand{\familydefault}{\sfdefault}

\usepackage{mathptmx} % skalierbare Formelschriften
\usepackage{tabto} % Tabulatoren
\usepackage{lastpage} % Letzte Seitenzahl auslesen
\usepackage{xstring} % Textmanipulation
\usepackage{setspace} % anpassbare Zeilenabstände

% Debugging:
%\usepackage{layout} % Layout-Informationen
%\usepackage{printlen} % Längenwerte ausgeben


\setlength{\PlakatFusszeileHoehe}{22mm}
\setlength{\PlakatFusszeilePositionUnten}{69mm}
\renewcommand{\PlakatFusszeileSpaltenzahl}{10}

\setlength{\PlakatSeiteRand}{46mm}
\setlength{\PlakatKopfzeileHoehe}{4cm}
\setlength{\PlakatKopfzeileAbstand}{46mm}
\setlength{\PlakatKopfzeilePositionUnten}{25mm}
\setlength{\PlakatUniversitaetLogoBreite}{87.2mm}
\setlength{\PlakatFakultaetLogoBreite}{55mm}
\setlength{\PlakatPositionLinks}{-2cm}
\setlength{\PlakatFakultaetslogoPositionKorrekturLinks}{3cm}
\setlength{\PlakatFakultaetslogoPositionKorrekturOben}{3mm}
\setlength{\PlakatFakultaetslogoAbstandRechts}{16mm}
\setlength{\PlakatUniversitaetslogoPositionKorrekturLinks}{-20mm}
\setlength{\PlakatSpaltenAbstand}{18mm}
\setlength{\PlakatSchriftNormalGroesse}{25pt}
\setlength{\PlakatSchriftNormalZeilenabstand}{32pt}
\setlength{\PlakatAufzaehlungAbstandLinks}{0.6em}

\newcommand{\PlakatTitelSchriftgroesse}{\fontsize{94}{113}\selectfont}

\newcommand{\PlakatTitelEins}[1]{{\PlakatTitelSchriftgroesse #1}\\[67pt]}
\newcommand{\PlakatTitelZwei}[1]{{\fontsize{67}{80}\selectfont #1}\\[22pt]}
\newcommand{\PlakatTitelDrei}[1]{{\fontsize{44}{53}\selectfont #1}\\[90pt]}

\newcommand{\PlakatKopfzeileSchriftgroesse}{\fontsize{42}{50}\selectfont}
\newcommand{\PlakatBildUnterschrift}[1]{{\\[-3mm]\fontsize{14}{18}\selectfont{}#1}}

\newcommand{\PlakatFusszeileSchrift}{\fontsize{12}{14}\selectfont}

\setlength{\PlakatSeiteHoehe}{62.8cm}
 % !!! NICHT ENTFERNEN !!!
\documentclass[a1, plainsections, landscape]{sciposter}
\newcommand{\PlakatFormat}{A1 quer}

\documentclass{scrlttr2} % Dokumentenklasse: KOMA-Skript Letter

% Anpassbare Kopf- und Fußzeilen für KOMA-Skript:
%\usepackage{scrlayer-scrpage}

\usepackage[utf8]{inputenc} % Textkodierung: UTF-8
\usepackage[T1]{fontenc} % Zeichensatzkodierung

\usepackage{calc} % Berechnungen

\usepackage[ngerman]{babel} % Deutsche Lokalisierung
\usepackage{graphicx} % Grafiken
\usepackage[absolute]{textpos} % Positionierung

% Silbentrennung:
\usepackage{hyphenat}
%\tolerance 2414
%\hbadness 2414
%\emergencystretch 1.5em
%\hfuzz 0.3pt
%\widowpenalty=10000     % Hurenkinder
%\clubpenalty=10000      % Schusterjungen
%\vfuzz \hfuzz

% Euro-Symbol:
\usepackage[gen]{eurosym}
\DeclareUnicodeCharacter{20AC}{\euro{}}

% Schriftart Helvetica:
\usepackage[scaled]{helvet}
\renewcommand{\familydefault}{\sfdefault}

\usepackage{mathptmx} % skalierbare Formelschriften
\usepackage{tabto} % Tabulatoren
\usepackage{lastpage} % Letzte Seitenzahl auslesen
\usepackage{xstring} % Textmanipulation
\usepackage{setspace} % anpassbare Zeilenabstände

% Debugging:
%\usepackage{layout} % Layout-Informationen
%\usepackage{printlen} % Längenwerte ausgeben


\setlength{\PlakatFusszeilePositionUnten}{47.5mm}
\renewcommand{\PlakatFusszeileSpaltenzahl}{6}

\setlength{\PlakatSeiteRand}{35mm}
\setlength{\PlakatKopfzeileHoehe}{4cm}
\setlength{\PlakatKopfzeileAbstand}{26mm}
\setlength{\PlakatKopfzeilePositionUnten}{25.5mm}
\setlength{\PlakatUniversitaetLogoBreite}{66.5mm}
\setlength{\PlakatFakultaetLogoBreite}{42mm}
\setlength{\PlakatPositionLinks}{-1cm}
\setlength{\PlakatFakultaetslogoPositionKorrekturLinks}{28.5mm}
\setlength{\PlakatFakultaetslogoPositionKorrekturOben}{2.5mm}
\setlength{\PlakatFakultaetslogoAbstandRechts}{12mm}
\setlength{\PlakatUniversitaetslogoPositionKorrekturLinks}{-9.5mm}
\setlength{\PlakatSpaltenAbstand}{14mm}
\setlength{\PlakatSchriftNormalGroesse}{22pt}
\setlength{\PlakatSchriftNormalZeilenabstand}{27pt}
\setlength{\PlakatFusszeileHoehe}{28mm}
\setlength{\PlakatAufzaehlungAbstandLinks}{0.7em}

\newcommand{\PlakatTitelSchriftgroesse}{\fontsize{67}{80}\selectfont}

\newcommand{\PlakatTitelEins}[1]{{\PlakatTitelSchriftgroesse #1}\\[38pt]}
\newcommand{\PlakatTitelZwei}[1]{{\fontsize{49}{53}\selectfont #1}\\[19pt]}
\newcommand{\PlakatTitelDrei}[1]{{\fontsize{39}{10}\selectfont #1}\\[68pt]}

\newcommand{\PlakatKopfzeileSchriftgroesse}{\fontsize{30}{39}}
\newcommand{\PlakatFusszeileSchrift}{\fontsize{11}{13.2}\selectfont}
\newcommand{\PlakatBildUnterschrift}[1]{{\\[-3mm]\fontsize{14}{18}\selectfont{}#1}}


\setlength{\PlakatSeiteHoehe}{42.9cm} % !!! NICHT ENTFERNEN !!!
%%%%%%%%%%%%%%%%%%%%%%%%%%%%%%%%%%%%%%%%%%%%%%%%%%%%%%%%%%%%%%%%%%%%%%%%%%%%%%%%
% TUM-Vorlage: Personenspezifische Informationen
%%%%%%%%%%%%%%%%%%%%%%%%%%%%%%%%%%%%%%%%%%%%%%%%%%%%%%%%%%%%%%%%%%%%%%%%%%%%%%%%
%
% Rechteinhaber:
%     Technische Universität München
%     https://www.tum.de
% 
% Gestaltung:
%     ediundsepp Gestaltungsgesellschaft, München
%     http://www.ediundsepp.de
% 
% Technische Umsetzung:
%     eWorks GmbH, Frankfurt am Main
%     http://www.eworks.de
%
%%%%%%%%%%%%%%%%%%%%%%%%%%%%%%%%%%%%%%%%%%%%%%%%%%%%%%%%%%%%%%%%%%%%%%%%%%%%%%%%

% Für die Person anpassen:

\newcommand{\PersonTitel}{Dr.~rer.~nat.}
\newcommand{\PersonVorname}{Erika}
\newcommand{\PersonNachname}{Mustermann}
\newcommand{\PersonStadt}{@Ort@}
\newcommand{\PersonAdresse}{%
    @Adresse@\\%
    @Plz@~\PersonStadt%
}
\newcommand{\PersonTelefon}{@Telefon@}
\newcommand{\PersonFax}{@Fax@}
\newcommand{\PersonEmail}{@E-Mail@}
\newcommand{\PersonWebseite}{@Web@}

\newcommand{\FakultaetAnsprechpartner}{@Ansprechpartner@}
% Fakultät:
\newcommand{\FakultaetName}{Informatics 15}
\newcommand{\LehrstuhlName}{Computer Graphics and Visualization}

\newcommand{\EinstellungBankName}{Bayerische Landesbank}
\newcommand{\EinstellungBankIBAN}{DE10700500000000024866}
\newcommand{\EinstellungBankBIC}{BYLADEMM}
\newcommand{\EinstellungSteuernummer}{143/241/80037}
\newcommand{\EinstellungUmsatzsteuerIdentifikationsnummer}{DE811193231}

\hyphenation{} % eigene Silbentrennung
                    % !!! DATEI ANPASSEN !!!
%%%%%%%%%%%%%%%%%%%%%%%%%%%%%%%%%%%%%%%%%%%%%%%%%%%%%%%%%%%%%%%%%%%%%%%%%%%%%%%%


\newcommand{\PlakatTitel}{Kickinect - VR Torwandschießen}


%%%%%%%%%%%%%%%%%%%%%%%%%%%%%%%%%%%%%%%%%%%%%%%%%%%%%%%%%%%%%%%%%%%%%%%%%%%%%%%%
%%%%%%%%%%%%%%%%%%%%%%%%%%%%%%%%%%%%%%%%%%%%%%%%%%%%%%%%%%%%%%%%%%%%%%%%%%%%%%%%
% EINSTELLUNGEN
%%%%%%%%%%%%%%%%%%%%%%%%%%%%%%%%%%%%%%%%%%%%%%%%%%%%%%%%%%%%%%%%%%%%%%%%%%%%%%%%

% Allgemein:
\newcommand{\AllgemeinGestalter}{ediundsepp Gestaltungsgesellschaft}
\newcommand{\AllgemeinErsteller}{eWorks GmbH}

% Universität:
\newcommand{\UniversitaetName}{Technische Universität München}
\newcommand{\UniversitaetAbkuerzung}{TUM}
\newcommand{\UniversitaetWebseite}{www.tum.de}
\newcommand{\UniversitaetLogoBreite}{19mm}
\newcommand{\UniversitaetLogoHoehe}{1cm}

\newcommand{\UniversitaetAdresse}{%
    Arcisstraße~21\\%
    80333~München%
}

\definecolor{UniversitaetFarbe}{RGB}{0,101,189}

% Seitenränder:
\newcommand{\Seitenrand}{1cm}
\newcommand{\SeitenrandLinks}{2cm}

% Text:
\newcommand{\TextOben}{24.2mm}
\newcommand{\TextSchriftgroesse}{11pt}
\newcommand{\TextZeilenabstand}{15pt}

% Aufzählungen:
\usepackage{enumitem}
\setlist[itemize]{topsep=0pt, partopsep=0pt, itemsep=0pt, parsep=0cm, leftmargin=5.2mm, itemindent=0cm} % weniger Platz vor und in Listen
\def\labelitemi{\raise0.2em\hbox{\fontsize{4pt}{4pt}\selectfont$\bullet$\hspace*{2.2mm}}}


%%%%%%%%%%%%%%%%%%%%%%%%%%%%%%%%%%%%%%%%%%%%%%%%%%%%%%%%%%%%%%%%%%%%%%%%%%%%%%%%
% DOKUMENT
%%%%%%%%%%%%%%%%%%%%%%%%%%%%%%%%%%%%%%%%%%%%%%%%%%%%%%%%%%%%%%%%%%%%%%%%%%%%%%%%

\usepackage[a4paper,
    top=\Seitenrand,
    bottom=\Seitenrand,
    inner=\SeitenrandLinks,
    outer=\Seitenrand,
    foot=0cm,
    head=0cm
]{geometry}

\textblockorigin{\SeitenrandLinks}{\Seitenrand} % Ursprung für Positionierung

% PDF-Einstellungen:
\usepackage[pdftex]{hyperref}
\hypersetup{
    hidelinks,
    pdfauthor={\PersonVorname{} \PersonNachname},
    pdftitle={\StellenausschreibungStelle},
    pdfproducer={\AllgemeinErsteller},
    pdfcreator={\AllgemeinGestalter}
}

\renewcommand*{\raggedsignature}{\raggedright}

\makeatletter
    \@setplength{refvpos}{\TextOben}
\makeatother

\KOMAoptions{
    fontsize=\TextSchriftgroesse,
    refline=nodate,
    foldmarks=false,
    firsthead=false,
    firstfoot=false,
    backaddress=no,
    addrfield=no,
    fromalign=false
}

\setkomavar{fromname}{\UniversitaetName}
\setkomavar{fromaddress}{\PersonAdresse}
\KOMAoption{fontsize}{10pt}

\newcommand{\StellenausschreibungHauptueberschrift}{{\fontsize{19}{25}\selectfont\vspace*{0.8mm}\StellenausschreibungStelle}}
\newcommand{\StellenausschreibungUnterueberschrift}[1]{\textbf{#1}\vspace*{-\baselineskip}}

\newcommand{\StellenausschreibungAdresseTrenner}{, \ignorespaces}
\newcommand{\StellenausschreibungAdresse}{{%
    \let\\\StellenausschreibungAdresseTrenner% Umdefinieren von "\\" zu ", "
    \PersonAdresse
}}

\AtBeginLetter{%
    % Dreizeiler:
    \begin{textblock*}{\textwidth}(0cm, 0cm)%
        \color{UniversitaetFarbe}%
        \fontsize{9}{11}\selectfont%
        \LehrstuhlName\\%
        \FakultaetName\\%
        \UniversitaetName%
    \end{textblock*}%
    \iftoggle{StellenausschreibungGrafikOpportunities}{%
        \begin{textblock*}{\textwidth}[1,0](\textwidth, 28.5mm-\Seitenrand)%
            \raggedleft%
            \includegraphics{./Ressourcen/Stellenausschreibung/Bilder/Opportunities.jpg}%
        \end{textblock*}%
    }{}
    % Logo:
    \begin{textblock*}{\UniversitaetLogoBreite}[1,0](\textwidth, 0cm)%
        \raggedleft%
        \includegraphics{./Ressourcen/_Bilder/Universitaet_Logo_RGB.pdf}%
    \end{textblock*}%
    \setlength{\baselineskip}{\TextZeilenabstand}%
    \setlength{\parindent}{0mm} % keine Einrückung am Absatzanfang
    \setlength{\parskip}{\baselineskip} % einzeiliger Abstand nach Absätzen
    \color{black}
}


\begin{document}
\raggedright
\begin{letter}{}
\opening{}

 % !!! NICHT ENTFERNEN !!!
%%%%%%%%%%%%%%%%%%%%%%%%%%%%%%%%%%%%%%%%%%%%%%%%%%%%%%%%%%%%%%%%%%%%%%%%%%%%%%%%


%%%%%%%%%%%%%%%%%%%%%%
%% 3-Spalten-Layout %%
%%%%%%%%%%%%%%%%%%%%%%

\PlakatKopfzeileLeer
\PlakatFusszeile{%
    \textbf{\UniversitaetName}\\
    \FakultaetName\\
    \LehrstuhlName%
}

\PlakatTitelEins{\PlakatTitel}
%\PlakatTitelZwei{Überschrift 2 läuft über gesamte Papierbreite}
\PlakatTitelDrei{Marcel Bruckner, Kevin Bein, }

\begin{multicols*}{3}

% Examples
%\PlakatUeberschrift{Schriftgröße 12 pt}

Dies ist die Vorlage für das Plakat (DIN \PlakatFormat{}) der Technischen
Universität München (TUM). Sie entspricht dem Corporate Design der TUM und ist für das Betriebssystem Windows getestet und mit TeX Live 2015 kompatibel.

Bitte geben Sie Ihren individuellen Text an den vorgesehenen Stellen ein. Nutzerhinweise finden Sie außerdem in der Liesmich.txt.

\ifx\PlakatBeschreibungKopfzeileUndAbsender\TRUE
\PlakatUeberschrift{Kopfzeile und Absender}

Es stehen mehrere Varianten zur Auswahl. Diese finden Sie auf den nächsten Seiten als Beispiel oder unter \href{http://\UniversitaetWebseite/cd}{\UniversitaetWebseite/cd}, welche für die verschiedenen Nutzergruppen vorgesehen sind.

\fi

\PlakatUeberschrift{Text}

Strukturieren Sie Ihr Plakat möglichst klar und beschränken Sie sich auf die wichtigsten Informationen. Entscheiden Sie sich für einen ein- oder mehrspaltigen Aufbau, linksbündig oder im Blocksatz. Und wählen Sie bei wenig Text eine entsprechend größere Schrift.

\PlakatUeberschrift{Bilder}

\ifx\PlakatBeschreibungKurz\TRUE % BEGINN kurze/lange Beschreibung

Auch wenn Sie Bilder integrieren möchten, können Sie innerhalb des vorgegebenen Rahmens variieren:

\begin{itemize}
\item Das Bild ist an die Breite der Textspalte angepasst
\item Das Bild geht über alle Textspalten
\item Das Bild ist seitenfüllend innerhalb des vorgesehenen Textrahmens (mit weißem Rand)
\item Das Bild ist seitenfüllend bis an den Papierrand (randabfallend) an mindestens drei Seiten
\end{itemize}

Bilder sollten immer eine Bildunterschrift tragen, die ggf. Informationen zum Bild und zur Autorenschaft enthält.
Das Bild sollte immer genügend Abstand zum Text haben. Gleiches gilt für Grafiken.

\else

Auch wenn Sie Bilder integrieren möchten, können Sie innerhalb des vorgegebenen Rahmens variieren:

\begin{itemize}
\item Das Bild ist an die Breite der Textspalte angepasst
\item Das Bild geht über alle Textspalten
\item Das Bild ist seitenfüllend innerhalb des vorgesehenen Textrahmens (mit weißem Rand)
\item Das Bild ist seitenfüllend bis an den Papierrand (randabfallend) an mindestens drei Seiten
\end{itemize}

Bilder sollten immer eine Bildunterschrift tragen, die ggf. Informationen zum Bild und zur Autorenschaft enthält. 
Das Bild sollte immer genügend Abstand zum Text haben. Gleiches gilt für Grafiken.

\ifx\PlakatBeschreibungBeispielbild\TRUE
\PlakatBild[0cm \PlakatBeschreibungBeispielbildBeschnitt{} 0cm 0cm]{./Ressourcen/_Bilder/Sternenhimmel.jpg}{Bildunterschrift, Autor etc}
\fi

\ifx\PlakatBeschreibungDruck\TRUE
\PlakatUeberschrift{Druck}

Achten Sie darauf, dass das Dokument in Originalgröße gedruckt wird, also keine Anpassung an die Druckränder bei Ihren Druckereinstellungen aktiviert ist.
Wählen Sie bitte unbedingt weiße Papiere statt Naturpapier mit starker Braun- oder Graufärbung. Denn nur sie entsprechen dem Charakter des Corporate Designs der TUM. Denken Sie daran, dass auch Papierstärke (Dicke) und Haptik einen Einfluss auf das Druckergebnis haben.
\fi

\fi % ENDE kurze/lange Beschreibung

Weitere Informationen zum Corporate Design der TUM finden Sie unter
\href{http://\UniversitaetWebseite/cd}{\UniversitaetWebseite/cd}


\PlakatUeberschrift{Abstract}

The goal of this project is to recreate the popular german game \textit{Torwandschießen} in VR. 

\PlakatUeberschrift{Game}

In \textit{Torwandschießen}, the player tries to kick a football through two openings in a wall.
One hole is aligned at the bottom left and the other hole is at the top right.
The player has five tries to kick the Ball through either of the two openings and collect points (1 point for the bottom and 3 points for the top hole).
Both holes are just slightly bigger than the diamter of the ball which makes it difficult to hit.


\PlakatUeberschrift{Implementation}

A Microsoft Kinect Sensor is used to track the player and his movements.
From this input data, a body skeleton is calculated and extracted (TODO: figure 2).
By using \textit{Linear Blend Skinning}, this skeleton is mapped onto the 3D Model of a football player.
For this game, it is especially important to detect movements of the lower body: the player's legs and foots.
Hitboxes are constructed around the players feet and the ball. 
If both collide in the virtual reality, we calculate and apply a velocity to the ball and accelarate it towards the Torwand.
When either of the two holes is hit perfectly, the corresponding points are accredited.
The game ends after five tries.

\end{multicols*}

\clearpage


%%%%%%%%%%%%%%%%%%%%%%
%% 5-Spalten-Layout %%
%%%%%%%%%%%%%%%%%%%%%%

%\PlakatKopfzeileMitDreizeiler
%\PlakatFusszeileLeer
%
%\PlakatTitelEins{\PlakatTitel}
%\PlakatTitelZwei{Überschrift 2 läuft über gesamte Papierbreite}
%\PlakatTitelDrei{Überschrift 3 läuft über gesamte Papierbreite oder gemäß Spaltenbreite}
%
%\begin{multicols*}{5}
%
%\setlength{\PlakatBeschreibungBeispielbildBeschnitt}{15cm} % Anpassen der Größe des Beispielbildes
%\PlakatUeberschrift{Schriftgröße 12 pt}

Dies ist die Vorlage für das Plakat (DIN \PlakatFormat{}) der Technischen
Universität München (TUM). Sie entspricht dem Corporate Design der TUM und ist für das Betriebssystem Windows getestet und mit TeX Live 2015 kompatibel.

Bitte geben Sie Ihren individuellen Text an den vorgesehenen Stellen ein. Nutzerhinweise finden Sie außerdem in der Liesmich.txt.

\ifx\PlakatBeschreibungKopfzeileUndAbsender\TRUE
\PlakatUeberschrift{Kopfzeile und Absender}

Es stehen mehrere Varianten zur Auswahl. Diese finden Sie auf den nächsten Seiten als Beispiel oder unter \href{http://\UniversitaetWebseite/cd}{\UniversitaetWebseite/cd}, welche für die verschiedenen Nutzergruppen vorgesehen sind.

\fi

\PlakatUeberschrift{Text}

Strukturieren Sie Ihr Plakat möglichst klar und beschränken Sie sich auf die wichtigsten Informationen. Entscheiden Sie sich für einen ein- oder mehrspaltigen Aufbau, linksbündig oder im Blocksatz. Und wählen Sie bei wenig Text eine entsprechend größere Schrift.

\PlakatUeberschrift{Bilder}

\ifx\PlakatBeschreibungKurz\TRUE % BEGINN kurze/lange Beschreibung

Auch wenn Sie Bilder integrieren möchten, können Sie innerhalb des vorgegebenen Rahmens variieren:

\begin{itemize}
\item Das Bild ist an die Breite der Textspalte angepasst
\item Das Bild geht über alle Textspalten
\item Das Bild ist seitenfüllend innerhalb des vorgesehenen Textrahmens (mit weißem Rand)
\item Das Bild ist seitenfüllend bis an den Papierrand (randabfallend) an mindestens drei Seiten
\end{itemize}

Bilder sollten immer eine Bildunterschrift tragen, die ggf. Informationen zum Bild und zur Autorenschaft enthält.
Das Bild sollte immer genügend Abstand zum Text haben. Gleiches gilt für Grafiken.

\else

Auch wenn Sie Bilder integrieren möchten, können Sie innerhalb des vorgegebenen Rahmens variieren:

\begin{itemize}
\item Das Bild ist an die Breite der Textspalte angepasst
\item Das Bild geht über alle Textspalten
\item Das Bild ist seitenfüllend innerhalb des vorgesehenen Textrahmens (mit weißem Rand)
\item Das Bild ist seitenfüllend bis an den Papierrand (randabfallend) an mindestens drei Seiten
\end{itemize}

Bilder sollten immer eine Bildunterschrift tragen, die ggf. Informationen zum Bild und zur Autorenschaft enthält. 
Das Bild sollte immer genügend Abstand zum Text haben. Gleiches gilt für Grafiken.

\ifx\PlakatBeschreibungBeispielbild\TRUE
\PlakatBild[0cm \PlakatBeschreibungBeispielbildBeschnitt{} 0cm 0cm]{./Ressourcen/_Bilder/Sternenhimmel.jpg}{Bildunterschrift, Autor etc}
\fi

\ifx\PlakatBeschreibungDruck\TRUE
\PlakatUeberschrift{Druck}

Achten Sie darauf, dass das Dokument in Originalgröße gedruckt wird, also keine Anpassung an die Druckränder bei Ihren Druckereinstellungen aktiviert ist.
Wählen Sie bitte unbedingt weiße Papiere statt Naturpapier mit starker Braun- oder Graufärbung. Denn nur sie entsprechen dem Charakter des Corporate Designs der TUM. Denken Sie daran, dass auch Papierstärke (Dicke) und Haptik einen Einfluss auf das Druckergebnis haben.
\fi

\fi % ENDE kurze/lange Beschreibung

Weitere Informationen zum Corporate Design der TUM finden Sie unter
\href{http://\UniversitaetWebseite/cd}{\UniversitaetWebseite/cd}

%
%
%\PlakatUeberschrift{Blindtext}
%
%\lipsum[1-8]
%
%
%\end{multicols*}
%
%\clearpage


%%%%%%%%%%%%%%%%%%%%%%
%% 1-Bild-Layout    %%
%%%%%%%%%%%%%%%%%%%%%%

%\PlakatKopfzeileMitEinzeiler
%\PlakatFusszeileMehrspaltig{
%    Hier kann ein längerer Text stehen, der in mehreren Spalten angeordnet wird.
%    \vfill\columnbreak
%    Durch \texttt{\textbackslash{}vfill\textbackslash{}columnbreak} lässt sich ein Spaltenwechsel erzwingen.
%    \vfill\columnbreak~
%    \vfill\columnbreak~
%    \vfill\columnbreak~
%    \vfill\columnbreak~
%    \vfill\columnbreak~
%    \vfill\columnbreak~
%    \vfill\columnbreak~
%}
%
%\PlakatTitelEins{\PlakatTitel}
%\PlakatTitelZwei{Überschrift 2 läuft über gesamte Papierbreite}
%
%\PlakatBildGanzseitig[0cm 18cm 0cm 0cm]{./Ressourcen/_Bilder/SternenhimmelQuer.jpg}{Bildunterschrift, Autor etc}
%
%
%\clearpage


%%%%%%%%%%%%%%%%%%%%%%%%%%%
%% Layout für wenig Text %%
%%%%%%%%%%%%%%%%%%%%%%%%%%%

%\PlakatKopfzeileMitEinzeiler
%\PlakatFusszeileLeer
%
%\PlakatTitelSchriftgroesse
%
%Wählen Sie für Plakate mit wenig Text gerne eine größere Schrift (siehe Auswahl für Überschriften). Achten Sie bitte darauf, dass die Schriftgröße Zweck und Gebrauch des Plakats entspricht.
%

%%%%%%%%%%%%%%%%%%%%%%%%%%%%%%%%%%%%%%%%%%%%%%%%%%%%%%%%%%%%%%%%%%%%%%%%%%%%%%%%
\end{document} % !!! NICHT ENTFERNEN !!!
%%%%%%%%%%%%%%%%%%%%%%%%%%%%%%%%%%%%%%%%%%%%%%%%%%%%%%%%%%%%%%%%%%%%%%%%%%%%%%%%
